\section*{Bab \ref{ch:g11:mutations} --- Mutasi dan Keragaman Genetik}

\begin{solution}{exo:g11:mutations:1}
Perubahan urutan \omterm{def:g10:universal-dna:nucleotide}{nukleotida} \omterm{def:g10:universal-dna:information}{DNA}, yang diwariskan kepada keturunan
\omterm{def:g10:cells-common-unit:cell}{selnya}. Substitusi, insersi, delesi.
\end{solution}

\begin{solution}{exo:g11:mutations:2}
\texttt{ATGCCAGAAT}: substitusi (T menjadi A pada kedudukan 6).
\texttt{ATGCCTGAT}: delesi satu A. \texttt{ATGCCTTGAAT}: insersi
satu T.
\end{solution}

\begin{solution}{exo:g11:mutations:3}
Somatik: dalam \omterm{def:g10:cells-common-unit:cell}{sel} tubuh, yang hanya diwariskan kepada keturunan \omterm{def:g10:cells-common-unit:cell}{sel}
itu di dalam individunya. \omterm{def:g10:cells-common-unit:cell}{Sel} kelamin: dalam gamet atau bakalnya, yang
diwariskan kepada keturunannya dan seluruh \omterm{def:g10:cells-common-unit:cell}{selnya}. Hanya \omterm{def:g11:mutations:mutation}{mutasi} \omterm{def:g10:cells-common-unit:cell}{sel}
kelamin yang masuk ke pewarisan sifat.
\end{solution}

\begin{solution}{exo:g11:mutations:4}
Cahaya ultraviolet, sinar-X, benzopiren asap tembakau (juga
keradioaktifan, aflatoksin). Ultraviolet melaskan dua timin yang
bertetangga pada satu untai, sehingga heliksnya terpuntir.
\end{solution}

\begin{solution}{exo:g11:mutations:5}
Untai komplementernya yang utuh: karena tiap untai menentukan untai
yang lain, urutan yang benar bagi ruas yang rusak dapat dibangun
kembali dari pasangannya.
\end{solution}

\begin{solution}{exo:g11:mutations:6}
Normal: sekitar 6\%; yang kurang perbaikannya: 0.01\%. Perbaikannya
meningkatkan ketahanan hidup sekitar 600 kali lipat pada dosis itu.
\end{solution}

\begin{solution}{exo:g11:mutations:7}
$20/10^9 = \num{2e-8}$. Tidak: cawannya hanya menyingkapkan \omterm{def:g10:cells-common-unit:cell}{sel} mana
yang sudah tahan; apakah antibiotiknya menyebabkan \omterm{def:g11:mutations:mutation}{mutasinya} menuntut
uji replikanya.
\end{solution}

\begin{solution}{exo:g11:mutations:8}
Tiap replikanya menerima \omterm{def:g10:cells-common-unit:cell}{sel} dari koloni asli yang sama. Jika
antibiotiknya menyebabkan ketahanan saat bersentuhan, \omterm{def:g10:cells-common-unit:cell}{sel} yang terkena
akan berupa segelintir \omterm{def:g10:cells-common-unit:cell}{sel} acak pada tiap replikanya, pada kedudukan
yang tak berhubungan. Menemukan ketahanannya selalu pada kedudukan yang
sama berarti koloni asli itu, yang tumbuh tanpa obatnya, sudah memuat
\omterm{def:g10:cells-common-unit:cell}{sel} yang tahan.
\end{solution}

\begin{solution}{exo:g11:mutations:9}
Kesalahan tiap salinan \omterm{def:g10:universal-dna:gene}{gennya}: $2 \times 3000 \times 10^{-9} =
\num{6e-6}$ (kedua untainya). Di antara $10^{12}$ \omterm{def:g10:cells-common-unit:cell}{sel} yang dihasilkan
pada pembelahan terakhirnya, sekitar $\num{6e6}$ membawa \omterm{def:g11:mutations:mutation}{mutasi} baru
pada \omterm{def:g10:universal-dna:gene}{gen} itu.
\end{solution}

\begin{solution}{exo:g11:mutations:10}
\omterm{def:g11:mutations:mutation}{Mutasi} somatik dalam sebuah \omterm{def:g10:cells-common-unit:cell}{sel} tunas memberi cabang itu warnanya;
seteknya adalah jaringan itu, sehingga seteknya menyimpannya. \omterm{def:g10:cells-common-unit:cell}{Sel}
kelamin bunganya berasal dari cabang yang sama, tetapi perubahan pada
satu \omterm{def:g10:universal-dna:gene}{alel} hanya diwariskan kepada separuh gametnya, dan bakal bijinya
dibuahi serbuk sari yang membawa \omterm{def:g10:universal-dna:gene}{alel} biasa; \omterm{def:g10:universal-dna:gene}{alel} yang baru itu, jika
resesif, tersembunyi pada kecambahnya.
\end{solution}

\begin{solution}{exo:g11:mutations:11}
\omterm{def:g10:cell-metabolism:metabolism}{Enzim} yang hilang itu memperbaiki kerusakan ultraviolet; ultraviolet
hanya menembus kulitnya. Ususnya tidak menerima ultraviolet, tidak
menderita lesi semacam itu, dan tidak memerlukan perbaikan itu.
\end{solution}

\begin{solution}{exo:g11:mutations:12}
Lesinya sebanding dengan dosisnya: sepersepuluh ultravioletnya berarti
sepersepuluh lesinya tiap \omterm{def:g10:cells-common-unit:cell}{sel} tiap jam --- tetap ribuan. Kulit terbakar
matahari adalah matinya banyak \omterm{def:g10:cells-common-unit:cell}{sel}; di bawah dosis itu \omterm{def:g10:cells-common-unit:cell}{selnya} bertahan,
tetapi lesi yang lolos dari perbaikan menjadi \omterm{def:g11:mutations:mutation}{mutasi}. Tidak terbakar
berarti kerusakannya masih dapat dilewati, bukan tidak ada.
\end{solution}

\begin{solution}{exo:g11:mutations:13}
Bukan keduanya dengan sendirinya. Di daerah bermalaria, satu salinannya
menaikkan ketahanan hidup, dan \omterm{def:g10:universal-dna:gene}{alelnya} tetap lazim meskipun penderita
dua salinannya sakit; di tempat malarianya tidak ada, ia hanya
menyebabkan penyakit dan menjadi langka. Lingkungannya, dan jumlah
salinannya, yang memutuskan nilai sebuah \omterm{def:g11:mutations:mutation}{mutasi}.
\end{solution}

\begin{solution}{exo:g11:mutations:14}
\omterm{def:g10:cells-common-unit:cell}{Sel} yang peka mati; \omterm{def:g10:cells-common-unit:cell}{sel} yang tahan, yang sudah hadir, membelah dengan
bebas dan dalam sehari biakannya seluruhnya tahan. Rangkaian yang
dihentikan lebih awal meninggalkan yang selamat terakhir --- yang
paling tahan --- tetap hidup untuk tumbuh kembali, dalam tubuh yang
kini diperkaya ketahanan; menghabiskan rangkaiannya membunuhnya sebelum
ia berlipat ganda.
\end{solution}

\begin{solution}{exo:g11:mutations:15}
Sebagian besar \omterm{def:g11:mutations:mutation}{mutasi} bersifat netral (di luar \omterm{def:g10:universal-dna:gene}{gen} atau tanpa pengaruh
pada \omterm{def:g10:chemistry-of-life:families}{proteinnya}), sebagian kecil merugikan, sangat sedikit yang
berguna; karena bersifat acak, \omterm{def:g11:mutations:mutation}{mutasi} terjadi tanpa memandang
kegunaannya, tetapi seleksi menyimpan yang berguna dan menyingkirkan
yang merugikan. Kegunaannya bergantung pada lingkungannya, sehingga
\omterm{def:g11:mutations:mutation}{mutasi} yang merugikan hari ini dapat menjadi \omterm{def:g11:mutations:mutation}{mutasi} yang menentukan
ketika keadaannya berubah. Sepanjang ribuan keturunan, dengan miliaran
individu, bahkan \omterm{def:g11:mutations:mutation}{mutasi} berguna yang langka pun muncul berulang kali:
evolusi dibangun di atas sisa yang terpilah, bukan di atas \omterm{def:g11:mutations:mutation}{mutasi}
rata-rata.
\end{solution}

\begin{solution}{pb:g11:mutations:1}
\textbf{1.} Agar tiap koloninya menjadi klon dengan asal yang
diketahui, dan agar koloninya cukup terpisah untuk dibedakan pada
replikanya.

\textbf{2.} Beberapa \omterm{def:g10:cells-common-unit:cell}{sel} dari puncak tiap koloninya, dalam susunan
geometri yang sama seperti pada cawannya, karena bantalannya ditekan
rata tanpa tergeser.

\textbf{3.} Bahwa \omterm{def:g10:cells-common-unit:cell}{sel} yang tahan hadir dalam kedua koloni asli itu:
tiga replika yang saling bebas tidak sepakat karena kebetulan.

\textbf{4.} Koloni tahan pada kedudukan acak yang tak berhubungan pada
ketiga cawannya, karena \omterm{def:g10:cells-common-unit:cell}{sel} mana yang terimbas akan berbeda dari
replika ke replika.

\textbf{5.} \omterm{def:g11:mutations:mutation}{Mutasi} sebelum sentuhannya: ketahanannya terpetakan pada
koloni asli, yang tidak pernah bertemu obatnya.

\textbf{6.} \omterm{def:g10:cells-common-unit:cell}{Sel} dari kedua kedudukan tahannya tumbuh pada
streptomisin; \omterm{def:g10:cells-common-unit:cell}{sel} dari kedudukan lainnya tidak.

\textbf{7.} Barangkali hanya sebagian: \omterm{def:g11:mutations:mutation}{mutasinya} terjadi pada suatu
pembelahan selama pertumbuhan koloninya, dan hanya keturunan \omterm{def:g10:cells-common-unit:cell}{sel} itu
yang tahan. Untuk memutuskannya, orang akan menanam sejumlah \omterm{def:g10:cells-common-unit:cell}{sel}
koloninya yang diketahui pada streptomisin dan menghitung yang selamat.

\textbf{8.} Pada pembelahan ke-20 dari 23: $1/2^{20-1}$\dots lebih
sederhana lagi, \omterm{def:g10:cells-common-unit:cell}{sel} mutannya lalu membelah 3 kali lagi: $2^3 = 8$ \omterm{def:g10:cells-common-unit:cell}{sel}
tahan dari $2^{23}$, sekitar satu tiap sejuta. Pada pembelahan ke-3:
satu dari 8 \omterm{def:g10:cells-common-unit:cell}{sel} yang lalu hadir, jadi seperdelapan koloninya.

\textbf{9.} Bahkan beberapa \omterm{def:g10:cells-common-unit:cell}{sel} tahan di antara sekitar $10^4$ \omterm{def:g10:cells-common-unit:cell}{sel} yang
diangkat beludrunya sudah cukup untuk mendirikan sebuah koloni pada
cawan antibiotiknya; ujinya mengenali kehadiran \omterm{def:g10:cells-common-unit:cell}{sel} yang tahan, bukan
bagiannya.

\textbf{10.} Alasannya bersandar pada kedudukan: hanya jika
geometrinya terjaga, sebuah koloni pada replikanya dapat disamakan
dengan sebuah koloni pada cawan aslinya.

\textbf{11.} $30/\num{10000} = \num{3e-3}$.

\textbf{12.} Sekitar $\num{3e-3}$ \omterm{def:g11:mutations:mutation}{mutasi} tiap $10^7$ pembelahan:
$\num{3e-10}$ tiap pembelahan.

\textbf{13.} Seorde dengan laju kesalahan pada satu \omterm{def:g10:universal-dna:nucleotide}{nukleotida}: \omterm{def:g11:mutations:mutation}{mutasi}
ketahanannya pada dasarnya adalah satu substitusi tertentu yang lolos
dari perbaikan.

\textbf{14.} $10^9$ pembelahan pada $\num{3e-10}$: rata-rata sepersekian
\omterm{def:g11:mutations:mutation}{mutasi}, sehingga sebagian besar biakannya memuat nol atau beberapa \omterm{def:g10:cells-common-unit:cell}{sel}
tahan (tiap \omterm{def:g11:mutations:mutation}{mutasinya} lalu dilipatgandakan pembelahan yang menyusulnya).

\textbf{15.} \omterm{def:g11:mutations:mutation}{Mutasi} pada pembelahan yang awal diwarisi bagian besar
biakannya, sehingga hitungannya besar; \omterm{def:g11:mutations:mutation}{mutasi} yang belakangan hanya
memberi beberapa \omterm{def:g10:cells-common-unit:cell}{sel} tahan. Karena saatnya acak, biakan yang diulang
memberi hitungan yang tersebar luas --- dan itu sendiri menjadi tanda
bahwa \omterm{def:g11:mutations:mutation}{mutasinya} muncul sebelum, bukan sebagai tanggapan atas, antibiotiknya.

\textbf{16.} Langka: \omterm{def:g11:mutations:mutation}{mutasi} jarang terjadi tiap pembelahan. Acak:
\omterm{def:g11:mutations:mutation}{mutasi} terjadi tanpa bergantung pada kegunaannya. Percobaannya
menegakkan keacakannya (alasan kedudukannya); hitungan Bagian III
mengukur kelangkaannya.

\textbf{17.} Antibiotiknya memilah populasinya: ia membunuh \omterm{def:g10:cells-common-unit:cell}{sel} yang
peka dan membiarkan \omterm{def:g10:cells-common-unit:cell}{sel} yang tahan berlipat ganda --- itulah seleksi.

\textbf{18.} \omterm{def:g11:mutations:mutation}{Mutasi} memasok \omterm{def:g10:universal-dna:gene}{alel} yang tahan, secara acak, sebelum ada
kebutuhan apa pun; antibiotiknya memilah populasinya, dan hanya
menyimpan pembawanya.

\textbf{19.} \omterm{def:g11:mutations:mutation}{Mutasinya} mengubah sebuah molekul yang terlibat dalam
kerja obatnya, bukan pertumbuhan atau rupa koloninya tanpa obatnya;
beda genetik punya pengaruh yang kasatmata hanya dalam lingkungan
tempat beda itu berarti.

\textbf{20.} Kedudukan yang serupa pada ketiga replikanya membuktikan
bahwa \omterm{def:g11:mutations:mutation}{mutasi} ketahanannya muncul sebelum, dan tanpa bergantung pada,
paparan antibiotiknya; laju yang terukur sekitar $\num{3e-10}$ tiap
pembelahan \omterm{def:g10:cells-common-unit:cell}{sel}.
\end{solution}
